\documentclass[11pt]{article}
\usepackage[a4paper,margin=28mm]{geometry}
\usepackage{amsmath,amssymb,amsthm,mathtools,booktabs,url}
\newtheorem{theorem}{Theorem}[section]
\newtheorem{lemma}[theorem]{Lemma}
\newtheorem{proposition}[theorem]{Proposition}
\newtheorem{corollary}[theorem]{Corollary}
\theoremstyle{remark}\newtheorem{remark}[theorem]{Remark}
\newcommand{\Z}{\mathbb Z}\newcommand{\Q}{\mathbb Q}
\newcommand{\C}{\mathbb C}\newcommand{\OO}{\mathcal O}
\DeclareMathOperator{\ord}{ord}\DeclareMathOperator{\diag}{diag}
\DeclareMathOperator{\rank}{rank}\DeclareMathOperator{\divisor}{div}
\DeclareMathOperator{\lcm}{lcm}
\title{Arithmetic determination of relations among lattice theta series\\
\large A formal Tate and finite-cover construction}
\author{Exploratory Proof Agent: Codex mathematical construction}
\date{22 September 2026}
\begin{document}\maketitle
\begin{abstract}
We describe an effective arithmetic reduction of relations among a finite family
of positive-definite integral lattice theta series to exact linear algebra.
An integral orthogonal sublattice expresses each series as a finite sum of
products of unary residue theta series. A formal product and norm calculation
identifies a common power of every unary factor with a section obtained from
Kato's algebraic theta unit and its initial term at the zero section. Normalization in a
common Kummer extension then produces regular sections of one line bundle,
and a divisor-degree argument supplies an explicit finite coefficient bound.
No Poisson summation, complex theta transformation, analytic valence formula,
or analytic Sturm bound enters the construction. We also prove exact stopping
criteria that often avoid the very large general bound, and give independent
examples with zero and nonzero theta kernels. An additional torsion-level
normalization, using exact projective theta ratios and four-torsion
coordinates, replaces the initial $K^s$ cover factor by a common
$2^{r-1}$ cover and yields a much smaller explicit coefficient bound.

The constructed sections are modular forms of weight $n/2$ in the
standard arithmetic sense on a common finite level of the algebraic
metaplectic moduli stack. We construct that level as an actual finite
etale moduli cover, identify the tautological half-Hodge line, prove
regularity at every cusp, and compute all prescribed Tate expansions.
Thus pivot columns give a modular-form basis of the theta image with
exact rational coordinates. No complex uniformization or analytic
comparison is used in these proofs. The finite level is not asserted
to be congruence, and no prescribed classical theta multiplier is claimed.
\end{abstract}

\section{Input and the finite-coefficient theorem}
Let $h\geq1$ and let $A_1,\ldots,A_h$ be symmetric positive-definite integral $n\times n$
matrices with even diagonal and $n>0$, and put
\[
 Q_i(x)=\tfrac12x^tA_ix,\qquad
 \theta_i(q)=\sum_{x\in\Z^n}q^{Q_i(x)}.
\]
Positive definiteness makes these series well-defined in $\Z[[q]]$.
For $n=0$ all input theta series are $1$, so the kernel is simply
$\{c:\sum_i c_i=0\}$ and the image basis is $\{1\}$.
The requested order problem is a special case: take the supplied complete
list of normalized representatives. The argument does not need an
unjustified implication from local elementary divisors to global class
completeness. Where representatives are supplied, their completeness remains
the stated input hypothesis.

For each $i$, the construction in Section~\ref{orthogonal} gives an integral
matrix $B_i$ of nonzero determinant such that
$B_i^tA_iB_i=\diag(m_{i1},\ldots,m_{in})$, with positive even $m_{ij}$.
Let $\mathcal U$ be the finite collection of distinct unary pairs $(m,\nu)$
occurring in the resulting coset decompositions, with $\nu$ taken modulo $m$.
Set
\begin{align*}
 N&=\lcm_{i,j}m_{ij},& L&=24N^2,& K&=(1+12N)^2-1,\\
 s&=|\mathcal U|,& I(L)&=L^3\prod_{p\mid L}(1-p^{-2}),&
 B_*&=\frac{nK^s I(L)}{48}.
\end{align*}
The integer $B_*$ is deliberately conservative.
Theorem~\ref{sharper-bound} proves the substantially smaller bound
$B_\dagger=\lfloor n2^{r-1}I(L)/(48L)\rfloor$, where $r$ counts
the distinct diagonal norms $m_{ij}$. The revised implementation uses
this latter bound; the first construction is retained as an independent
arithmetic route and for comparison.

\begin{theorem}[Arithmetic finite-coefficient determination]\label{main}
Using the algebraic Tate uniformization and algebraic theta-unit theorems
specified in Section~\ref{tate}, the construction below gives
\[
 \ker_{\C}\bigl(c\longmapsto\textstyle\sum_i c_i\theta_i\bigr)
 =\ker_{\C} M_{B_*},\qquad
 M_B=\bigl([q^a]\theta_i\bigr)_{0\leq a\leq B,\,1\leq i\leq h}.
\]
In particular the kernel is the complex span of an effectively computable
rational basis. Pivot columns give a basis of the formal theta image and
exact rational coordinates for every theta series in that basis.
The same series are the exact Tate expansions of standard arithmetic
modular forms of weight $n/2$ on one common finite metaplectic level,
regular at every cusp. Consequently the pivot theta series form a
linearly independent basis of modular forms for the image. The actual
level and its half-Hodge line are established in
Corollary~\ref{arithmetic-modularity}.
\end{theorem}

We use the standard arithmetic definition of half-integral-weight
modular forms. On a finite representable etale level of the algebraic
metaplectic moduli problem, such a form is a section of a tensor power
of its tautological half-Hodge line, regular also at the cusps of the
compactification. The metaplectic object and its finite-level meaning
are specified in \cite[Definition 3.1.1 and Remark 3.1.2]{candelori}.
We give the test-object description and verify the finite-level map in
Corollary~\ref{arithmetic-modularity}, so the word ``modular'' does not
mean an arbitrary line bundle endowed with a formal expansion.
This permits the usual noncongruence finite-level scope; it does not
claim a specified Shimura congruence level or multiplier. All proofs
below take place in arithmetic moduli and formal completions. The
relation to analytic terminology is discussed separately in
Section~\ref{scope} and is not used in the construction.

\section{An orthogonal sublattice and unary factors}\label{orthogonal}
\begin{lemma}\label{orthogonal-lemma}
For an even integral positive-definite matrix $A$, there is an integral
$B$ with nonzero determinant and $B^tAB=\diag(m_1,\ldots,m_n)$, where
each $m_j$ is positive and even. For
$H=B^{-1}\Z^n/\Z^n$, choose representatives $t$. Then
$\nu_j=m_jt_j$ is integral, and
\[
 \theta_A(q)=\sum_{t\in H}\prod_{j=1}^n U_{m_j,\nu_j}(q),\qquad
 U_{m,\nu}(q)=\sum_{r\equiv\nu\, (m)}q^{r^2/(2m)}.
\]
All terms of each product have integral total exponent.
\end{lemma}
\begin{proof}
Rational Gram--Schmidt produces an orthogonal basis over $\Q$. Clearing
denominators of its columns produces $B$. Its diagonal norms are even
because its columns are integral vectors in the even lattice. The index of
$B\Z^n$ in $\Z^n$ is $|\det B|$.

If $x=Bt\in\Z^n$, then
$\diag(m_j)t=B^tAx\in\Z^n$. Hence $m_jt_j\in\Z$.
The coset of $x$ consists precisely of $B(t+k)$ for $k\in\Z^n$, and
\[
 Q(B(t+k))=\sum_j\frac{m_j}{2}(t_j+k_j)^2.
\]
Putting $r_j=m_j(t_j+k_j)$ gives the asserted product. Since $B(t+k)$
is integral, its total exponent is integral. Each bounded exponent involves
only finitely many vectors, so all regroupings are formal and legitimate.
\end{proof}

For example, for $A=\left(\begin{smallmatrix}2&1\\1&12\end{smallmatrix}\right)$,
take $B=\left(\begin{smallmatrix}1&-1\\0&2\end{smallmatrix}\right)$.
Then $B^tAB=\diag(2,46)$ and
\[
 \theta_{x^2+xy+6y^2}=U_{2,0}U_{46,0}+U_{2,1}U_{46,23}.
\]
This reduction works without a parity restriction on the rank.

\section{Formal products and algebraic theta units}\label{tate}
We use Kato's algebraic existence and norm characterization
\cite[Proposition 1.3(1)]{kato}: for $c>1$ prime to $6$ there is a
base-change compatible rational function ${}_c\vartheta_E$ with divisor
$c^2[0]-[E[c]]$, invariant under the norm for multiplication by every
$d$ prime to $c$. Its algebraic construction uses divisors and Picard
functoriality; see also the modern statement in
\cite[Theorem 4.1.1]{klz}. Kato's original complex proof of the product formula is
\emph{not} used here.

For the formal quotient we use the actual positive Tate datum
$R=k[[Q]]$, with $k$ of characteristic zero, the torus $\mathbb G_m$,
and the period lattice $Q^{\Z}$. Its polarization is the symmetric
pairing $(Q^r,Q^s)\mapsto Q^{rs}$; for $r\ne0$, $Q^{r^2}$ lies in
the ideal of definition. Thus the completeness, normality and positive
polarization hypotheses of Mumford's construction
\cite[Theorem 3.10, Corollary 4.9, Theorem 4.10]{mumford} are satisfied.
We use the explicit annular model and proper algebraization of
\cite[Corollary 2.2.4 and Section 2.5]{conrad}. Its generic fiber is a
geometrically connected elliptic curve, with torsion and multiplication
represented by $z^d=Q^i$ and $z\mapsto z^d$. No assertion about an
unspecified period datum over an arbitrary discrete valuation ring is needed.

\begin{lemma}[Constant comparison on a proper formal model]\label{constant}
Let $R=k[[t]]$ have characteristic zero and let $\mathcal E/R$ be a
proper semistable model with geometrically connected smooth generic
fiber. Choose a finite affine formal covering
$\widehat{\mathcal E}=\bigcup_\alpha\operatorname{Spf}(A_\alpha)$
by noetherian normal domains $A_\alpha$ that are $t$-torsion-free.
Suppose compatible elements
$u_\alpha\in\operatorname{Frac}(A_\alpha[1/t])$ have no pole at any
height-one prime of $A_\alpha[1/t]$. Then they are the restrictions
of a constant in $k((t))$.
\end{lemma}
\begin{proof}
The height-one intersection criterion for a noetherian normal domain
implies $u_\alpha\in A_\alpha[1/t]$. Choose $b_\alpha\geq0$ with
$t^{b_\alpha}u_\alpha\in A_\alpha$. The covering is finite, so a
single $b=\max_\alpha b_\alpha$ works on all charts. Compatibility
then gives a global regular function on $\widehat{\mathcal E}$.
The theorem on formal functions identifies this with an element of
$H^0(\mathcal E,\mathcal O_{\mathcal E})$: this module is finite over
the complete ring $R$, hence is already complete. Its restriction to
the proper geometrically connected smooth generic fiber is constant.
Dividing by $t^b$ proves the assertion.
\end{proof}


Define
\[
 \gamma_Q(z)=\prod_{r\geq0}(1-Q^rz)\prod_{r\geq1}(1-Q^rz^{-1}),
 \qquad \gamma_Q(Qz)=-z^{-1}\gamma_Q(z).
\]
For $K_c=c^2-1$ and $p_c=-c(c-1)/2$, put
\[
 f_c(z)=Q^{K_c/12}(-z)^{p_c}
              \frac{\gamma_Q(z)^{c^2}}{\gamma_Q(z^c)}.
\]

We explain why this lemma applies to the comparison used below.
For the fixed integer $c>1$, after a finite scalar extension containing
$\mu_c$ and the ramified extension $Q=t^r$, $r=2c$, use the actual
$r$-gon Tate model of \cite[Section 2.5]{conrad}. Its finitely many
annular charts have rings
\[
 A_j=k[[t]]\langle x_j,y_j\rangle/(x_jy_j-t),\qquad 0\leq j<r.
\]
On the infinite chain before taking the period quotient, the
multiplicative coordinate is $z=t^jx_j$, so
$z^{-1}=t^{-j-1}y_j$. The quotient identifies the last overlap by
$z\mapsto Qz$. These rings are normal; if necessary refine the cover
by completions of algebraic affine charts of the same proper model.

On the $j$th chart the two products in $\gamma_Q(z)$ are
\[
 \prod_{\ell\geq0}(1-t^{r\ell+j}x_j),\qquad
 \prod_{\ell\geq1}(1-t^{r\ell-j-1}y_j).
\]
Every exponent is nonnegative. Apart from $1-x_j$ when $j=0$ and
$1-y_j$ when $j=r-1$, all factors form a unit of $A_j$.
For $\gamma_Q(z^c)$ the corresponding exponents are
$r\ell+cj$ and $E=r\ell-c(j+1)$. Only finitely many $E$ are nonpositive.
Since $c\mid r$, a negative $E$ satisfies $E\leq-c$, and
\[
 1-t^Ey_j^c=-t^Ey_j^c(1-t^{-E-c}x_j^c).
\]
The monomial is a unit after inverting $t$; the last factor is a unit
of $A_j$ unless $E=-c$, when it is $1-x_j^c$.
The case $E=0$ gives $1-y_j^c$. All infinite remainders are units.
Consequently $\gamma_Q(z^c)\in A_j[1/t]$, and $f_c$ lies in its
fraction field. Its only horizontal divisor factors are the simple
linear factors of $1-x_j^c$ or $1-y_j^c$ and the indicated factors
of $\gamma_Q(z)$.

Indeed the $c$-torsion representatives are
$z=\zeta_c^b t^{ra/c}$, $0\leq a,b<c$. Since $c\mid r$, these extend
to smooth sections of the $r$-gon; the factors just listed are their
local equations. Hence the horizontal divisor of $f_c$ is exactly
$c^2[0]-[E_Q[c]]$. Kato's algebraic function has the same divisor.
On the same normal algebraic model its other divisor components are
vertical. After passing to the compatible formal charts and inverting
$t$, no vertical component remains; units remain units, and the split
smooth torsion sections retain their displayed local equations.
Thus the ratio of Kato's function and $f_c$ has order zero at every
height-one prime of every $A_j[1/t]$.

Both functions are invariant under $z\mapsto Qz$, so their local
ratios agree under the actual quotient overlap maps. Lemma~\ref{constant}
therefore makes their ratio a constant over the finite extended
coefficient field. The norm computations for $d=2,3$ force this
constant to satisfy $u^3=u^8=1$, hence $u=1$. The resulting identity
holds already over the original coefficient field. The auxiliary
extension $Q=t^{2c}$ is only used in this comparison proof and does
not enlarge the later modular level or its degree bound.

\begin{proposition}[A formal norm calculation]\label{norm}
If $c>1$ and $(c,6)=1$, then $f_c$ is the full Tate expansion of
${}_c\vartheta_{E_Q}$.
\end{proposition}
\begin{proof}
The translation identity for $\gamma$ proves periodicity of $f_c$.
Its divisor is $c^2[0]-[E[c]]$. Let $(d,c)=1$. The $d^2$ preimages of $z$
are $z_{ij}=\zeta_d^jQ^{i/d}z^{1/d}$ for $0\leq i,j<d$.
The finite identity $\prod_j(1-\zeta_d^jw)=1-w^d$ yields
\[
 \prod_{i,j}\gamma_Q(z_{ij})=\gamma_Q(z).
\]
Write $ci=dk_i+r_i$ with $0\leq r_i<d$, and set
\begin{align*}
 S&=\sum_i k_i=\frac{(c-1)(d-1)}2,\\
 E&=\sum_i\left(r_i k_i+\frac d2 k_i(k_i-1)\right)
 =\frac{(c^2-1)(d-1)(2d-1)}{12}
       -\frac{d(c-1)(d-1)}4.
\end{align*}
The latter equality follows by expanding $r_i=ci-dk_i$ and using that
the $r_i$ permute $0,\ldots,d-1$. Translation gives
\[
 \prod_{i,j}\gamma_Q(z_{ij}^{c})
 =(-1)^{-S}z^{-cS}Q^{-E}\gamma_Q(z^c).
\]
Also $\prod_{i,j}(-z_{ij})=(-1)^{d^2}z^dQ^{d(d-1)/2}$.
The norm ratio therefore has possible factors with exponents
\[
 (d^2-1)p_c+S,\quad (d-1)p_c+cS,\quad
 (d^2-1)K_c/12+p_c d(d-1)/2+E.
\]
The last two are zero. The first is
$(c-1)(d-1)(1-c(d+1))/2$, an even integer: check $d$ odd and even
separately. Thus $\operatorname{Norm}_{[d]}f_c=f_c$.

The ratio of Kato's function to $f_c$ has no divisor and is a constant
$u$ by Lemma~\ref{constant}. Norm invariance gives $u^{d^2}=u$.
Taking $d=2,3$ gives $u^3=u^8=1$, whence $u=1$.
\end{proof}

Let $B_2(T)=T^2-T+1/6$ and define the formal Siegel product
\[
 g_{a,b}(Q)=-Q^{B_2(a)/2}e^{\pi i b(a-1)}
                         \gamma_Q(Q^a e^{2\pi i b}),\quad 0\leq a<1.
\]
The roots of unity in this notation are algebraic constants. Put
$\eta(Q)=Q^{1/24}\prod_{r\geq1}(1-Q^r)$.
The formal Jacobi triple product gives, for $-m/2\leq\nu<m/2$,
\[
 U_{m,\nu}(q)=-e^{-\pi i(a-1)/2}\eta(q^m)g_{a,1/2}(q^m),
 \qquad a=\nu/m+1/2.
\]
The leading exponent is fixed by
$1/24+B_2(a)/2=(a-1/2)^2/2$; there is no unspecified formal unit.
For completeness, an arithmetic proof of the triple product is as follows.
Write $(x;t)_r=\prod_{0\leq j<r}(1-xt^j)$. Euler coefficient recursions give
\[
 (-z;t)_\infty=\sum_{r\geq0}\frac{t^{r(r-1)/2}z^r}{(t;t)_r}.
\]
The product $F(z)=(-z;t)_\infty(-t/z;t)_\infty(t;t)_\infty$ satisfies
$F(tz)=z^{-1}F(z)$. If $F=\sum a_rz^r$, then
$a_r=a_0t^{r(r-1)/2}$. Its constant coefficient is
$(t;t)_\infty\sum_{r\geq0}t^{r^2}/(t;t)_r^2=1$: the sum is the
partition generating function, decomposed by the size of its Durfee square.
This proves the product identity coefficientwise, with normalization $a_0=1$.
Substitution $t=Q,z=Q^a$ gives the stated unary identity; bounded exponents
still involve finitely many terms, including when $a=0$.

Take $c=1+12N$ and $K=c^2-1$. Evaluating Proposition~\ref{norm} at
$z=Q^ae^{2\pi ib}$, with $Na,Nb\in\Z$ and
$(a,b)\notin\Z^2$, gives
${}_c\vartheta(P_{a,b})=g_{a,b}^K$.
Indeed the difference of phase exponents is
\[
 (c-1)(a-c/2)+(c-1)b[1-(c-1)a]
 =12A-6Nc+12B-144AB\in2\Z,
\]
where $a=A/N,b=B/N$. Since $24N\mid K$, the remaining triple-product
phase also vanishes. Consequently
\begin{equation}\label{power}
 U_{m,\nu}(q)^K=\mathcal D_2(q^m)^{K/48}
       {}_c\vartheta_{E_{q^m}}(P_{\nu/m+1/2,\,1/2}).
\end{equation}

Here $\mathcal D_2(Q)=Q^2\prod_{r\geq1}(1-Q^r)^{48}$.
The exponent $K/48$ is integral: $K=24N(1+6N)$ and $N$ is even,
because it is the least common multiple of a nonempty family of positive
even integers $m_{ij}$. The torsion points actually used above have
$b=1/2$ and are therefore nonzero.
Thus \eqref{power} is a purely formal product identity so far.

\begin{lemma}[An algebraic discriminant square from a theta unit]\label{d2}
On the moduli of elliptic curves in characteristic zero there is a
canonical nowhere-vanishing section $D_2$ of $\omega^{24}$ with exact
Tate expansion $\mathcal D_2(Q)\,(dz/z)^{24}$.
\end{lemma}
\begin{proof}
For $c=5$, Kato's function has a zero of order $24$ along the identity
section and no other zero or pole in a neighborhood of it. If $\mathcal I$
is the ideal of the identity, then
$\mathcal I^{24}/\mathcal I^{25}\simeq\omega^{24}$.
Define $D_2=5\operatorname{in}_0({}_5\vartheta)$ under this canonical
identification. The exact divisor makes this initial term everywhere
nonzero, and base-change compatibility makes it canonical.

Put $z=1+t$ and $P=\prod_{r\geq1}(1-Q^r)$. Then
$\gamma_Q(1+t)=-tP^2+O(t^2)$ and
$\gamma_Q((1+t)^5)=-5tP^2+O(t^2)$.
Proposition~\ref{norm}, for which $K_5=24$ and $p_5=-10$, gives
\[
 {}_5\vartheta(1+t)=\tfrac15 Q^2P^{48}t^{24}+O(t^{25}).
\]
The class of $t$ is $dz/z$ at the identity. Multiplying the initial term
by $5$ proves the full expansion, including its constant factor.
\end{proof}

\section{Algebraic sections and every cusp}
Work over a number field containing the required roots of unity and take
a geometric component $X$ of the full-level modular curve $X(L)$,
$L=24N^2$, containing the designated Tate cusp. The ambient fine moduli
scheme is smooth and projective in characteristic zero: see
\cite[Theorems 4.1.1 and 4.2.1(2), Lemma 4.2.4]{conrad}.
The selected geometric component and cusp descend to a finite number
field extension. Its completed cusp rings have $q=t^L$
\cite[Lemma 4.1.7 and (4.2.1)]{conrad}. The Hodge line on the
compactification is the invariant cotangent line at the identity of the
smooth locus of the universal generalized elliptic curve. Its cusp frame
is $d\log z$.
Its universal source elliptic curve has a designated cyclic
isogeny $\phi_m:E\to E_m$ of degree $m$, with target $N$-torsion
sections. At the designated cusp this is
$E_q\to E_{q^m}$, $z\mapsto z^m$
\cite[Example 2.5.1]{conrad}.
The level suffices because $mN\mid L$. More explicitly, a target point
$P=\zeta_N^B q^{mA/N}$, with integer coordinates $A,B$, has source lift
$\widetilde P=\zeta_{mN}^B q^{A/N}$, an $mN$-torsion point.
These integers are distinct from the rational characteristics $a,b$ used
in the formal product identities.
The chosen full-level structure supplies this lift and the kernel
$\mu_m$ of the quotient isogeny; consequently the selected target
$N$-torsion sections exist on this algebraic modular curve.

Let $\omega$ be the Hodge line of the source. In \eqref{power}, replace
the formal $\mathcal D_2$ by the algebraic $D_2$ of Lemma~\ref{d2},
pull back the target section and multiply by $m^{-K/2}$. This normalization is
essential: $\phi_m^*(dz'/z')=m\,dz/z$ at the designated cusp.
It defines a section $F_{m,\nu}$ of $\omega^{K/2}$ whose full expansion
in the source Tate differential is exactly $U_{m,\nu}^K$.

For clarity, the calculation at each cusp is an independent base
change of the same global algebraic theta unit. If $Q_x$ is the target
Tate period at a cusp $x$, Proposition~\ref{norm} is applied there with
its formal variable $Q$ replaced by $Q_x$. Base-change naturality of
Kato's function and of the initial-term construction of $D_2$ puts both
factors of the target section in this same completed field. No formula
identifying the completed parameters at two different cusps is used.

At another cusp the same nonzero torsion point has Tate coordinates
$(a',b')$ with $0\leq a'<1$. If $a'=0$, then $b'\notin\Z$.
Thus the leading coefficient of its Siegel product is nonzero, and
the order of the target section is
\[
 K/24+K B_2(a')/2=K(a'-1/2)^2/2\geq0.
\]
To justify the local Hodge comparison, write the intersection of the
cyclic kernel with the torus as $\mu_e$, $e\mid m$. Quotienting by
$\mu_e$ gives $z\mapsto z^e$; the remaining component quotient enlarges
the period lattice and is the identity on torus coordinates. These are
actual formal Tate maps \cite[Section 2.5]{conrad}, after a finite
parameter extension if necessary. Their differential factor is the
nonzero constant $e$ in characteristic zero. Pullback multiplies valuations by positive
ramification indices and does not add a pole from the Hodge comparison.
Hence $F_{m,\nu}$ is regular at every cusp. It is nonzero and nowhere
vanishing on the open modular curve: $D_2$ is a unit there, and the
nonzero $N$-torsion point is not $c$-torsion.

For the full level problem with a fixed Weil pairing, finite
level-structure counting gives at most $I(L)/(2L)$ cusps, each of width
$L$. Primitive pairs modulo $L$, up to simultaneous sign, give this
count; equivalently divide the $I(L)$ symplectic markings by the cusp
stabilizer of order $2L$. The cusps of the selected component form a
subset of this finite set. Let $C$ denote their reduced divisor.
Naturality of $D_2$ and its exact Tate product prove
\[
 \divisor(D_2)=2LC,\qquad
 24\deg\omega=2L\deg C\leq I(L).
\]
Only this upper bound is needed. We do not consume the proof of
geometric connectedness of the entire fixed-pairing modular curve in
\cite[Theorem 4.2.1(1)]{conrad}, which uses complex connectedness.

\section{A common root line and an effective jet bound}
Set $D=(L/12)C$ and $M=\omega\otimes\OO_X(-D)$. Since $24\mid L$,
$D/2$ is an integral divisor. Write $1_D$ for the canonical rational
section of $\OO_X(D)$, whose divisor is $D$. Then
\[
 d_M=D_2/1_D^{24}\in H^0(X,M^{24})
\]
is nowhere zero. The line-bundle root torsor
\[
 T=\{v\in M:v^{24}=d_M\}\longrightarrow X
\]
is finite etale of degree $24$: locally its equation is $z^{24}=a$
with $a$ a unit, and its derivative $24z^{23}$ is invertible.
On $T$, put $\varepsilon=v\,1_D$. It has divisor the pullback of $D$,
so it supplies an isomorphism
\[
 H_T=\OO_T(p^*D/2),\qquad
 H_T^2\simeq p^*\omega,\qquad 1_{H_T}^2\longmapsto\varepsilon.
\]
The equality of the two divisors makes this a regular isomorphism at
every point, not just a rational comparison.

Here the common Tate frame can be specified explicitly. Write
$q=t^L$, $E(t)=\prod_{r\geq1}(1-t^{Lr})$ and $\delta=d\log z$.
The regular frames $e_D=t^{-L/12}1_D$ and $e_M=\delta/e_D$ satisfy
$d_M=E(t)^{48}e_M^{24}$. Choose the actual formal root
$v=E(t)^2e_M$ and the geometric component $X_0$ containing this branch.
The map $p_0:X_0\to X$ has degree at most $24$ and is unramified
at the selected cusp. For $H=H_T|_{X_0}$ the frame
\[
 h_T=E(t)^{-1}t^{-L/24}1_H
\]
satisfies $h_T^2=\delta$. This construction avoids any parity
assumption on the degree of an individual component's Hodge line,
and uses no analytic connectedness theorem or eta transformation law.

\begin{lemma}[Simultaneous global roots]\label{roots}
Let $X$ be a smooth projective curve in characteristic zero, $H$ a line
bundle and $F_1,\ldots,F_s$ nonzero sections of $H^K$.
Specify roots of these sections in one completed algebraic closure at a
point, using the same local frame. There exists a finite map $\pi:Y\to X$
of smooth projective curves of degree at most $K^s$, on which these roots
are regular sections of $\pi^*H$ with the specified expansions.
\end{lemma}
\begin{proof}
Choose a rational section $h$ of $H$ and write $F_j=f_jh^K$.
In the specified completed algebraic closure choose roots $r_j^K=f_j$
and normalize $X$ in the field generated by all $r_j$. This is one field
component, not an arbitrary collection of independently chosen components.
We perform this construction over the algebraic closure of the constant
field and keep the component containing the prescribed Tate embedding.
Thus $Y$ is geometrically connected. Its degree is at most $K^s$; all
finitely many chosen algebraic data can subsequently descend to a finite
constant extension. The rational sections $s_j=r_j\pi^*h$
satisfy $s_j^K=\pi^*F_j$. At every point $y$,
$K\ord_y(s_j)=\ord_y(\pi^*F_j)\geq0$, proving regularity.
The completed embedding gives the requested expansions.
\end{proof}

Apply Lemma~\ref{roots} over $X_0$ to the pulled-back sections
$p_0^*F_j$ and the just constructed common frame $h_T$.
Prescribe $r_j=U_j(h_T/h)$ simultaneously. All these roots belong to
the same Tate completion: $2m\mid L$ implies $U_j(t^L)\in\overline\Q[[t]]$. Therefore $s_j=U_jh_T$ at this cusp.
Every product of $n$ unary factors lies in the same line
$\mathcal L=(\pi^*H)^n$. By Lemma~\ref{orthogonal-lemma}, the theta
series are the expansions of genuine sections $\sigma_i$ of this line,
in the common frame $h_T^n$.

For the composite map $Y\to X_0\to X$ the degree is at most $24K^s$.
Consequently
\[
 \deg\mathcal L\leq\frac n2\,24K^s\deg\omega
       \leq\frac{nK^sI(L)}2.
\]
At the chosen cusp, $w=\ord_y(q)\geq L$. Since $L\geq24$ and
$B_*=nK^sI(L)/48$, a linear combination vanishing through $q^{B_*}$
has section order at least
\[
 w(B_*+1)\geq L(B_*+1)
 >\frac{nK^sI(L)}2\geq\deg\mathcal L.
\]
A nonzero global section of a line bundle on a smooth proper curve has
zero divisor of degree equal to the degree of that line. The section must
therefore be zero. Expanding in the common frame proves that the full
formal linear combination is zero. This proves Theorem~\ref{main}.

Notice that the numerical bound needs only $n,N,s$ and the finite group
order $I(L)$. The kernel algorithm need not compute the root torsor,
factor a Kummer extension, find its optimal degree, or enumerate cusp
inertia cycles in order to use this sufficient bound.

\section{Exact algorithms and practical stopping certificates}
Each coefficient up to a given $B$ is computable by exact ellipsoid
enumeration. A rational $LDL^t$ decomposition supplies positive rational
diagonal entries and permits recursive completion of squares. At each
stage integer interval endpoints are obtained with integer square roots
and exact rational inequalities. The search is finite, enumerates every
vector with $Q(x)\leq B$, and counts both signs without approximation.
Exact Gaussian elimination over $\Q$ then computes $\ker M_B$, pivot
columns and the coordinates of all columns. Base extension to $\C$
preserves the matrix rank and kernel dimension.

\begin{lemma}[Earlier certificates]\label{early}
For every $B\geq0$, $\ker\Theta\subseteq\ker M_B$.
If $M_B$ has full column rank, the complete theta kernel is zero.
More generally, if a space $W$ of rigorously established full-series
relations has dimension $h-\rank M_B$, then
$W=\ker\Theta=\ker M_B$.
\end{lemma}
\begin{proof}
Truncation of a full relation proves the inclusion. In the second case
$W\subseteq\ker\Theta\subseteq\ker M_B$ and the two outer spaces have
the same dimension.
\end{proof}

In particular, a singleton theta family has complete kernel zero already
at $B=0$, since its constant term is $1$. This assertion is independent of
modularity. An integral basis-change certificate
$P^tAP=A'$, $\det P=\pm1$, proves a full-series relation by a norm-preserving
bijection; merely matching a finite prefix proves no such relation.

The companion implementation \texttt{exact-engine/theta\_exact.py} returns
full kernel and coordinate claims only when these stopping certificates
justify them. Deficient truncations without sufficient full-series
certificates remain explicitly unresolved. The companion \texttt{solve\_universal.py} derives the bound internally
and computes prefixes at $0,1,2,4,\ldots$, capped at $B_\dagger$. It stops at
an earlier certificate when available; otherwise it computes through
$B_\dagger$ and applies Theorem~\ref{sharper-bound}. This is a terminating algorithm
on every valid finite input and accepts no externally asserted bound.
Its default command only inspects the bound; \texttt{--complete} executes
the algorithm without an attempt budget. The symbolic bound and actual
completed coefficient enumeration are recorded separately. The theorem
proves termination, not practical tractability of every input.

\section{Two arithmetic examples}
\subsection{A complete order with three independent classes}
In $(\Q^2,X^2+20Y^2)$, use bases
\[
 T_1=I,\qquad T_2=\begin{pmatrix}3&1\\0&1\end{pmatrix},\qquad
 T_3=\diag(2,1/2).
\]
Their normalized forms are
$x^2+20y^2$, $3x^2+2xy+7y^2$, and $4x^2+5y^2$.
The unnormalized restriction for $T_2$ is \emph{three times} the second
form; thus no false unscaled rational isometry is asserted.
All normalized Gram matrices have Smith factors $(2,40)$.

For a reduced form $ax^2+bxy+cy^2$ with Gram determinant $80$,
$|b|\leq a\leq c$ implies $3a^2\leq80$, hence $a\leq5$.
Enumerating these finitely many possibilities and imposing the Smith
condition leaves $(1,0,20),(3,\pm2,7),(4,0,5)$. The two signs are
integrally equivalent. The minima $1,3,4$ distinguish the three classes,
and all three have been realized in the fixed rational space. Their
coefficient rows at $0,1,3$ form
\[
 \begin{pmatrix}1&1&1\\2&0&0\\0&2&0\end{pmatrix},
 \qquad\det=4.
\]
The full kernel is zero. The companion appendix proves the completeness,
content normalization and rational-similarity assertions in detail.

\subsection{A full nonzero relation between distinct classes}
Let $D_n=\{x\in\Z^n:\sum x_i\equiv0\pmod2\}$ and
$D_n^+=D_n\cup((1/2,\ldots,1/2)+D_n)$ for $8\mid n$.
These are even unimodular lattices. Write
\[
 a=\sum_{r\in\Z}u^{4r^2},\quad
 b=\sum_{r\in\Z}(-1)^ru^{4r^2},\quad
 c=\sum_{r\in\Z}u^{(2r+1)^2}.
\]
The four-dimensional Hadamard matrix divided by $2$ bijects integral
vectors of odd coordinate sum with one parity half of the half-integral
vectors, preserving norm. Reflection in one coordinate exchanges the
two half-integral parity halves. Consequently $a^4=b^4+c^4$ as a
formal identity. Parity selection gives the norm series of $D_n^+$ as
$(a^n+b^n+c^n)/2$.
The polynomial identity
\[
 \left(\frac{X^2+Y^2+Z^2}{2}\right)^2
       =\frac{X^4+Y^4+Z^4}{2}\quad(X=Y+Z)
\]
then proves $\theta_{D_8^+\perp D_8^+}=\theta_{D_{16}^+}$ coefficientwise.
Substitution $q=u^8$ is injective, so this is the original integral
theta identity, not a finite check.

The root nonorthogonality graph of the orthogonal sum is disconnected,
whereas that of $D_{16}^+$ is connected. The two lattices are not
isometric. Since their determinants are both $1$, any lattice similarity
would have positive scale $s$ satisfying $s^{16}=1$, hence would be an
isometry. They are distinct rational-similarity classes in the same
order. For these two supplied columns the full kernel is exactly
$\C(1,-1)$. No unproved classification of the entire rank-sixteen order
is required for this example.

\section{Arithmetic normalization on torsion-level moduli}\label{refinement}
The following refinements give full formal expansions on specified
arithmetic level structures. In particular they remove a common-unit
ambiguity in a projective theta comparison. The comparison of the resulting
half-weight lines is kept separate from these identities.

\subsection{A theta fourth power from two-torsion coordinates}
Write $P(Q)=\prod_{r\geq1}(1-Q^r)$. In the standard Tate Weierstrass
model with invariant differential $\delta=dz/z$, the $x$-coordinate is
\[
 X(z;Q)=\sum_{r\in\Z}\frac{Q^rz}{(1-Q^rz)^2}
       -2\sum_{r\geq1}\frac{rQ^r}{1-Q^r}.
\]
For negative indices the summands are rewritten using
$v/(1-v)^2=v^{-1}/(1-v^{-1})^2$; the resulting expressions are formal
on each Tate annulus. To verify that these are actual algebraic coordinates,
use the proper polygon model of \cite[Section 2.5]{conrad}, after a finite
parameter extension. On its annular charts the displayed inversion rewrites
all negative-index terms with nonnegative parameter order. Away from the
identity all denominators are units; at the identity only the indicated
double pole remains. Thus $X$ is a section of
$\OO(2[0])$ on the formal model. If a further chart refinement introduces
vertical denominators, a single parameter power clears them, since the
cover is finite. Proper formal functions for this invertible sheaf
algebraize that section, just as in Lemma~\ref{constant}.

Periodicity follows by shifting the summation index, and
$X(z^{-1})=X(z)$. The function $Y=\tfrac12z\,\partial X/\partial z$
similarly algebraizes with a triple pole at the identity and is odd under
inversion. The degree-two function $X$ gives the quotient by inversion,
so $Y^2$ is a cubic in $X$. Its leading coefficient is one by the Laurent
terms $X(1+s)=s^{-2}+O(s^{-1})$ and
$Y(1+s)=-s^{-3}+O(s^{-2})$. Its three roots are the $X$-coordinates
of the nonzero two-torsion points. This proves
$Y^2=\prod_j(X-e_j)$ and $dX/(2Y)=dz/z$ on the actual algebraic
Tate curve, rather than assuming that an arbitrary periodic formal
expression is rational.

\begin{lemma}[Exact basic theta normalization]\label{basic-fourth}
For $\vartheta(q)=\sum_{r\in\Z}q^{r^2}$ one has the full formal identity
\[
 \vartheta(q)^4=4\bigl(X(q;q^2)-X(-1;q^2)\bigr).
\]
\end{lemma}
\begin{proof}
First the multiplicative addition formula is
\begin{equation}\label{tate-difference}
 X(z;Q)-X(w;Q)
 =\frac{-wP(Q)^4\gamma_Q(z/w)\gamma_Q(zw)}
        {\gamma_Q(z)^2\gamma_Q(w)^2}.
\end{equation}
For generic $w$, both sides are periodic in $z$, have zero divisor
$[w]+[w^{-1}]$ and pole divisor $2[1]$. Their quotient is constant by
Lemma~\ref{constant}, applied to a proper Tate model after a finite
extension splitting the points. The local horizontal factors of each
$\gamma$ are its indicated simple zeros; the remaining products are
units, as in the proof of Proposition~\ref{norm}. The same comparison
may be made with $w$ indeterminate and then specialized as a rational
identity. At $z=1+s$, use
$\gamma_Q(1+s)=-sP(Q)^2+O(s^2)$ and
$\gamma_Q(w^{-1})=-w^{-1}\gamma_Q(w)$. The right side has leading term
$s^{-2}$, equal to that of $X(1+s)-X(w)$, so the constant is one.
Specialization to nonzero two-torsion is valid; the two simple zeros
then coalesce to a double zero.

Set $Q=q^2,z=q,w=-1$, and abbreviate
\[
 A=\prod_{r\geq1}(1-q^{2r-1}),\quad
 B=\prod_{r\geq1}(1+q^{2r-1}),\quad
 C=\prod_{r\geq1}(1+q^{2r}).
\]
Then $\gamma_{q^2}(q)=A^2$, $\gamma_{q^2}(-q)=B^2$ and
$\gamma_{q^2}(-1)=2C^2$. The Euler identity
\[
 \prod_{r\geq1}(1-T^{2r-1})\prod_{r\geq1}(1+T^r)=1
\]
gives $ABC=1$. Thus \eqref{tate-difference} becomes
\[
 4(X(q;q^2)-X(-1;q^2))
   =P(q^2)^4 B^4/(A^4C^4)=P(q^2)^4B^8.
\]
The formal triple product already proved gives
$\vartheta(q)=P(q^2)B^2$, proving the identity in all degrees.
\end{proof}

\begin{proposition}[A square root supplied by four-torsion]\label{level-four-root}
On a characteristic-zero elliptic curve with labelled nonzero two-torsion
$P_j=(e_j,0)$, write
\[
 Y^2=(X-e_1)(X-e_2)(X-e_3),\qquad\delta=dX/(2Y).
\]
Choose a four-torsion section $R$ with $2R=P_2$. Then
\[
 g=\frac{(X(R)-e_1)^2-(e_1-e_2)(e_1-e_3)}{2Y(R)},
 \qquad \alpha=2g\delta
\]
defines a weight-one section on this arithmetic level structure, and
$\alpha^2=4(e_2-e_1)\delta^2$. At the Tate marking
$Q=q^2$, $P_1:z=-1$, $P_2:z=q$, $R:z=q^{1/2}$, its expansion is
\[
 \alpha=\vartheta(q)^2\delta.
\]
\end{proposition}
\begin{proof}
Put $u=X-e_1$, $a=e_1-e_2$, $b=e_1-e_3$. The tangent-and-chord law
on $Y^2=u^3+(a+b)u^2+ab u$ gives
\[
 X(2R)-e_1=\frac{(u(R)^2-ab)^2}{4Y(R)^2}.
\]
Since $2R=P_2$, this proves $g^2=e_2-e_1$. The denominator is nonzero
because $R$ has exact order four. Under a Weierstrass coordinate change
$g$ has weight one; hence $2g\delta$ is intrinsic. At the specified
Tate marking, $X(R)=q^{1/2}+O(q)$,
$Y(R)=\tfrac12q^{1/2}+O(q)$, $e_1=-1/4+O(q^2)$ and
$e_2=2q+O(q^2)$, so $g=1/2+O(q^{1/2})$.
Lemma~\ref{basic-fourth} identifies its square, and the leading sign
selects exactly $\vartheta(q)^2$.
\end{proof}

At a full-level cusp, the nonzero two-torsion representatives are
$-1,Q^{1/2},-Q^{1/2}$, in some order. The difference of the last two
$X$-coordinates has $Q$-order $1/2$; either difference involving the
first has order zero. Thus $\alpha$ extends at every cusp, with
$Q$-order zero or $1/4$, and it is nowhere zero on the open curve.
On a full level $M$ divisible by eight, $Q=t^M$, so
$\divisor(\alpha)=2D$ for an integral effective cusp divisor $D$.
Consequently $H=\OO(D)$, $s=1_D$ and $s^2\mapsto\alpha$ give a
specific algebraic square root $H^2\simeq\omega$. Its prescribed Tate
frame is $h=\vartheta(q)^{-1}s$, with $h^2=\delta$ and
$s=\vartheta(q)h$. Here $q=Q^{1/2}$ on the specified cusp.

For $\vartheta(q^d)$, pull this construction back along the cyclic
isogeny $E_q\to E_{q^{2d}}$, $z\mapsto z^{2d}$. The target Hodge
section must be multiplied by $(2d)^{-1}$, since the target differential
pulls back to $2d\,d\log z$. A full source level divisible by $2dM$
supplies the required target torsion lifts. This produces a normalized
half-weight section for each $d$, but does not yet assert that all such
half-weight lines have been identified on one torsion level.

\subsection{The exact finite theta-multiplier character}
The remaining fourth-root phase of Proposition~\ref{level-four-root}
can also be computed arithmetically. This identifies the basic
multiplier, without asserting compatibility for different isogenies.

\begin{proposition}[A full-four-level character]\label{four-character}
Let $b=(R_1,R_2)$ be an ordered full-four basis, $P_i=2R_i$,
$P_3=P_1+P_2$, and $\zeta=e_4(R_1,R_2)$. For the right action on bases
of the subgroup
\[
 G=\{g\in\mathrm{SL}_2(\Z/4):g\bmod2\text{ preserves }xy\},
\]
the section $\alpha$ of Proposition~\ref{level-four-root} satisfies
$\alpha(bg)=\chi(g)\alpha(b)$, where
\[
 \chi(S)=\zeta^{-1},\qquad\chi(T^2)=1,\qquad
 S=\begin{pmatrix}0&-1\\1&0\end{pmatrix},\quad
 T^2=\begin{pmatrix}1&2\\0&1\end{pmatrix}.
\]
This is exactly the genus-one even discriminant character
$\lambda$ of \cite[Section 2.2]{candelori-key}, with these Weil-pairing
and ordered-basis conventions.
\end{proposition}
\begin{proof}
Translate $e_1$ to zero and write
$Y^2=X(X+a)(X+b)$, $v=X(R_1)$, $u=X(R_2)$ and $y_i=Y(R_i)$.
The duplication conditions imply
\[
 v^2=ab,\qquad u^2+2au+ab=0,\qquad
 y_1^2=v(v+a)(v+b),\quad y_2^2=u(u+a)(u+b).
\]
Our root is $g_{12}=-a(u+b)/y_2$. Under $S$ the ordered basis becomes
$(R_2,-R_1)$, and direct substitution into its defining formula gives
\[
 g_{12}(bS)=-a(v+b)/y_1,\qquad
 \chi(S)=\frac{(v+b)y_2}{(u+b)y_1}.
\]
To fix this fourth-root phase, use normalized Miller functions.
If $\ell_i$ is the tangent linear function at $R_i$, then
\[
 \divisor(\ell_i)=2[R_i]+[P_i]-3[0],\qquad
 f_i=\ell_i^2/(X-e_i),\quad\divisor(f_i)=4[R_i]-4[0].
\]
At the identity, in parameter $-X/Y$, $f_i$ has leading coefficient
one. Algebraic Weil reciprocity, applied after moving the supports
and then taking the limit at the identity, gives
$e_4(R_1,R_2)=f_1(R_2)/f_2(R_1)$; the general normalization sign is
$(-1)^4=1$. The tangents pass through $P_i$, so
\[
 \ell_1(R_2)=y_2-y_1u/v,\qquad
 \ell_2(R_1)=y_1-y_2(v+a)/(u+a).
\]
For $t=y_1/y_2$, substitution yields
\[
 e_4(R_1,R_2)
 =\frac{(v+a)(u+a)^2(v-tu)^2}
 {uv^2((u+a)t-v-a)^2}
 =\frac{(u+b)t}{v+b}=\chi(S)^{-1}.
\]
For explicit verification, after clearing denominators the middle
identity reads
\[
 (v+a)(v+b)(u+a)^2(v-tu)^2
 -uv^2(u+b)t((u+a)t-v-a)^2=0.
\]
It follows by substituting $v^2=ab$, $u^2=-2au-ab$ and
$t^2u(u+a)(u+b)=v(v+a)(v+b)$. The denominators do not vanish on
the full-four-level moduli problem: the distinct halves have order
four and the tangent divisors avoid the evaluation points.
This proves $\chi(S)=\zeta^{-1}$.

Adding $P_1$ sends $u$ to $ab/u$ and $y_2$ to $-ab y_2/u^2$;
the formula $(u^2-ab)/(2y_2)$ is unchanged. Hence $\chi(T^2)=1$.
The mod-two kernel in $G$ consists of $I+2A$ with
$\operatorname{tr}(A)=0$ and has order eight. It is generated by
$S^2=-I$, $T^2$ and $ST^2S^{-1}$; the remaining mod-two image is
$\{I,S\}$. Thus $|G|=16$ and $S,T^2$ generate it. The computed
ratios, independent of the curve, multiply under successive changes
of marking, and therefore define the claimed full character.
The defining values of $\lambda$ in the cited source are
$\lambda(S)=-1$, $\lambda(T^2)=0$ in additive $\Z/4$ notation.
They agree with the values above under $j\mapsto\zeta^j$.
\end{proof}

This also fixes the descent direction. The theta multiplier line
$M_\Theta$ uses the torsor of isomorphisms from $E[4]$ to the standard
module, whereas our ordered basis is an isomorphism in the opposite
direction. Its level frame consequently satisfies
$m_{bg}=\lambda(g)^{-1}m_b$. Therefore
\[
 \alpha(b)\otimes m_b\in\omega\otimes M_\Theta
\]
is independent of the full-four marking and descends to the moduli
problem with the even characteristic $P_3$.
For the rigidified line $\Theta=\OO(P_3)$, let $J=\pi_*\Theta$.
Its identity-evaluation section $e\in J^{-1}$ is nowhere zero on
the open moduli problem, because $P_3$ is disjoint from the identity.
Thus the rule
\[
 J^{-2}\xrightarrow{\ \sim\ }\omega\otimes M_\Theta,
 \qquad e^2\longmapsto\alpha(b)\otimes m_b
\]
is an explicit algebraic key-formula isomorphism for this genus-one
even characteristic. Its full Tate normalization is fixed by
Proposition~\ref{level-four-root}; no claim that all units on an open
modular curve are constant has been used.

\subsection{Projective theta values without an unknown formal unit}
\begin{proposition}[Exact projective unary ratios]\label{projective-ratios}
For every positive even $m$ there are algebraic rational functions
$R_{m,\nu}$ on an arithmetic full $2m$-level curve, with a specified
compatible theta marking, such that their complete Tate expansions are
\[
 R_{m,\nu}(q)=\frac{U_{m,\nu}(q)}{U_{m,0}(q)},
 \qquad U_{m,0}(q)=\vartheta(q^{m/2}).
\]
\end{proposition}
\begin{proof}
Use the rigidified totally symmetric degree-$m$ line
$\mathcal L_m=\OO_E(m[0])\otimes\omega_E^m$.
The algebraic full-level domination of symmetric theta structures in
\cite[Remark 2.6.5]{candelori} supplies the structure on full level $2m$.
The algebraic Schr\"odinger comparison
\cite[Theorems 4.1.1 and 4.1.3]{candelori} identifies
$\pi_*\mathcal L_m$ with the standard $m$-dimensional representation
tensored with one line. Evaluation at the identity therefore has
components $s_{m,\nu}$ in a common line. Their quotients are rational
functions, independent of the trivialization of that line.

Here is the complete formal normalization, including relative phases.
On the Tate multiplicative cover the section equation is
\[
 f(qz)=q^{-m/2}z^{-m}f(z).
\]
Indeed $z^{-m/2}\gamma_q(z)^m$ has this multiplier and divisor $m[0]$;
the rigidification changes only the common base frame. On full $2m$
level choose the symmetric theta marking with operators
\[
 (Af)(z)=f(\zeta_m z),\qquad
 (Bf)(z)=q^{1/(2m)}z f(q^{1/m}z).
\]
They satisfy $A^m=B^m=1$, $AB=\zeta_m BA$, and inversion conjugates
each to its inverse. The extra root $q^{1/(2m)}$ belongs to the specified
full level, not to an independently chosen completion extension.
To justify this marking globally, the theta structure supplied by
\cite[Remark 2.6.5]{candelori} induces the prescribed full $m$-level
structure. Its difference from the displayed symmetric lifts, which
have the same underlying level, belongs to $K(m)[2]$ by
\cite[Lemma 2.5.3]{candelori}. On full level this is a constant finite
group. On the chosen connected Tate branch the difference is therefore
one constant element; composing with that same global theta-group
automorphism aligns the lifts. This correction changes only finite
signs and adds neither a parameter-dependent factor nor a new cover.

The exact eigenbasis is
\[
 \vartheta_{m,\nu}(z)=\sum_{r\equiv\nu\, (m)}q^{r^2/(2m)}z^r.
\]
The positive quadratic exponent makes each Laurent sum convergent
in the formal sense on every annulus: substituting $z=t^jx$ leaves
parameter orders tending quadratically to infinity in both directions.
The prescribed automorphy glues these as formal sections of the actual
line $\mathcal L_m$. A finite cover bounds any vertical denominators;
proper formal functions for that line then algebraize the sections.
The generic section space has dimension $m$ by Riemann--Roch on the
elliptic curve. Thus these are actual section vectors, not only
solutions in an unspecified Laurent space.
Its $A$-eigenvalue is $\zeta_m^\nu$, and
$B\vartheta_{m,\nu}=\vartheta_{m,\nu+1}$ without an extra factor.
For any section $\sum a_rz^r$, its equation implies
$a_{r+m}=q^{r+m/2}a_r$. Thus each $A$-eigenspace is one-dimensional.
An intertwiner sends the standard basis to
$b_\nu\vartheta_{m,\nu}$; the $B$ action forces all $b_\nu$ to be
one and the same nonzero scalar. Evaluating at $z=1$ and dividing by
the zeroth component cancels this scalar and the rigidification frame.
The denominator is not identically zero, since $U_{m,0}$ has constant
term one. This proves the asserted rational function and every
coefficient of its expansion. No half-Hodge comparison has been used.
\end{proof}

These two constructions reduce an arithmetic half-weight comparison to
basic theta lines: formally,
$U_{m,\nu}=R_{m,\nu}\vartheta(q^{m/2})$.
Once a candidate section has been put in a specified half-Hodge line,
its regularity need not be guessed from its rational-function factors.
Its $K$th power equals $F_{m,\nu}$ by the injection of the function
field into the designated Tate completion, since both complete
expansions agree. Thus its order at every point is nonnegative by the
already proved regularity of $F_{m,\nu}$.

\subsection{A common arithmetic line and a smaller coefficient bound}
The torsion-level constructions above give a much smaller common cover
than adjoining every $K$th root separately. This observation closes the
common-line construction for the arithmetic coefficient theorem; it does
not identify the resulting finite cover as a congruence cover.

\begin{theorem}[A smaller arithmetic finite-coefficient bound]\label{sharper-bound}
Let $r$ be the number of distinct positive even integers $m_{ij}$ in the
orthogonal decompositions, and retain $N,L$ as above. Then one may replace
$B_*$ in Theorem~\ref{main} by
\[
 B_\dagger=
 \left\lfloor\frac{n\,2^{r-1}I(L)}{48L}\right\rfloor.
\]
The proof uses the arithmetic level constructions of this section and
none of the complex comparison in Section~\ref{scope}.
\end{theorem}
\begin{proof}
For each distinct $m$, the full source level $L=24N^2$ supplies the cyclic
$m$-isogeny $\phi_m:E\to E_m$ and a full-four marking on its target:
$4m\mid L$. At the distinguished cusp the target period is $q^m$.
Pull back the section $\alpha$ of Proposition~\ref{level-four-root}
and divide by $m$, using the Hodge pullback isomorphism. Call the
resulting source weight-one section $a_m$. Its exact expansion is
\[
 a_m=\vartheta(q^{m/2})^2\delta=U_{m,0}(q)^2\delta.
\]
This normalization fixes the full series and its sign simultaneously.

At an arbitrary source cusp write $q_x=t^L$. If the toric part of the
cyclic kernel is $\mu_e$, the residual component part has order
$b=m/e$. The target period satisfies
$Q_x^b=\zeta q_x^e$, so
\[
 \ord_t(Q_x)=Le^2/m.
\]
The Hodge map there is multiplication by $e$; the factor $e/m$ is a
nonzero constant and introduces no zero or pole. The earlier
four-torsion calculation shows that $a_m$ has order either zero or
$Le^2/(4m)$. Since $8m\mid L$, both possibilities are even integers.
It has no interior zero, since the two distinct two-torsion coordinates
never coincide. Hence $\divisor(a_m)=2D_m$ for an effective integral
cusp divisor $D_m$ on $X$. Set
\[
 H_m=\OO_X(D_m),\qquad e_m=1_{D_m},\qquad
 H_m^2\xrightarrow{\sim}\omega,\quad e_m^2\longmapsto a_m.
\]
The equality of divisors proves that this is an everywhere invertible
line-bundle map, including the boundary.

Fix one value $m_0$ and write $H_0=H_{m_0}$. For each of the other
$r-1$ values, $T_m=H_m\otimes H_0^{-1}$ has a specified trivialization
$T_m^2\simeq\OO_X$. Its unit-frame torsor is therefore finite etale
of degree two over the entire compact curve. Form their fiber product
and choose the component with the following common Tate branch:
\[
 t_m=\frac{U_{m_0,0}(q)}{U_{m,0}(q)}\,
             \frac{e_m}{e_{m_0}},\qquad t_m^2=1.
\]
All these series have constant term one, and both displayed sections
are nonvanishing at the distinguished cusp. Thus the branch already
exists over its completed local ring. This gives a finite etale map
$p:Y\to X$ of degree at most $2^{r-1}$, unramified at that cusp.
For $r=1$ take $Y=X$.

Dividing $p^*e_m$ by the tautological frame $t_m$ gives a section
$s_{m,0}$ of the common line $p^*H_0$. With
$h_0=U_{m_0,0}^{-1}e_{m_0}$, one has $h_0^2=\delta$ and
$s_{m,0}=U_{m,0}h_0$ at the selected cusp. Define the rational section
\[
 s_{m,\nu}=p^*R_{m,\nu}\,s_{m,0}.
\]
Proposition~\ref{projective-ratios} gives its exact expansion
$U_{m,\nu}h_0$. Under $(p^*H_0)^K\simeq p^*\omega^{K/2}$ its
$K$th power consequently equals $p^*F_{m,\nu}$: their full
expansions agree, and the function field injects into the cusp
completion. Since the latter is regular everywhere,
$K\ord_y(s_{m,\nu})\geq0$ at every point $y$, so $s_{m,\nu}$ is
regular. This also handles possible poles in the rational factors $R$.

The original coset expressions now give all $\theta_i$ as sections of
one line $\mathcal L=(p^*H_0)^n$, with a single common Tate frame
$h_0^n$. Using $24\deg\omega\leq I(L)$ yields
\[
 \deg\mathcal L=\frac n2\deg(p)\deg\omega
       \leq\frac{n\,2^{r-1}I(L)}{48}.
\]
The selected cusp still has $\ord(q)=L$. A nonzero section cannot
vanish to order $L(B_\dagger+1)$, since this exceeds the displayed
degree bound by the definition of the floor. Thus vanishing of the
coefficients through $B_\dagger$ forces the complete section, and
therefore its formal theta series, to vanish. Base change to $\C$
gives the same statement for all complex linear combinations.
\end{proof}

The new bound replaces the factor $K^s$ by $2^{r-1}/L$. Its numerical
advantage is independent of the early stopping certificates of
Lemma~\ref{early}: those certificates can still finish long
before even $B_\dagger$. Neither the compact etale torsors used here
nor their characters have been asserted to descend to a principal
congruence level merely because they are finite.

\subsection{Arithmetic modularity and the image basis}
\begin{corollary}[Modularity on a standard finite metaplectic level]
\label{arithmetic-modularity}
The sections constructed in Theorem~\ref{sharper-bound} are standard
arithmetic modular forms of weight $n/2$ at a common finite metaplectic
level. Their complete Tate expansions are the input $\theta_i$.
Pivot columns of $M_{B_\dagger}$ give a linearly independent
modular-form basis of their image, and exact rational elimination gives
the coordinates of every $\theta_i$ in that basis.
\end{corollary}
\begin{proof}
Let $X^\circ$ be the open full-level curve and
$\mathcal X_{1/2}=\sqrt[2]{\omega/X^\circ}$ the algebraic metaplectic
stack of \cite[Definition 3.1.1]{candelori}. Its object over a test
scheme $T$ is a level-marked elliptic curve, a line $R$, and a specified
isomorphism $R^2\simeq\omega_T$. Denote its tautological line by
$\lambda$. The actual isomorphism
$H^2\simeq p^*\omega$, with $H=p^*H_0$, defines
\[
 j:Y^\circ\longrightarrow\mathcal X_{1/2},\qquad j^*\lambda=H.
\]
This is a genuine finite metaplectic level: for a test object
$(T,R,R^2\simeq\omega_T)$, the fiber product
$Y^\circ\times_{\mathcal X_{1/2}}T$ is the torsor
\[
 \operatorname{Isom}_{\mathrm{sq}}(R,H)
       \longrightarrow Y^\circ\times_{X^\circ}T
\]
of isomorphisms compatible with the two specified square maps.
It is a representable finite etale $\mu_2$-torsor. Since $p$ is finite
etale, $j$ is representable finite etale, of degree $2\deg p$.
Thus both the level and its weight line arise from the pre-existing
metaplectic moduli functor.

The same square isomorphism holds on the entire compact curve $Y$.
The sections of $H^n$ are regular at every cusp by the proof of
Theorem~\ref{sharper-bound}, and their expansions in the common
nonvanishing frame $h_0^n$ are exactly $\theta_i(q)$. Hence these are
modular forms of weight $n/2$ in the arithmetic definition, with
holomorphy at the boundary, not merely sections on the open curve.
Their linear relations coincide with their formal-series relations
by injectivity at the chosen cusp. Theorem~\ref{sharper-bound} and
rational elimination now give the asserted basis and coordinates.
\end{proof}

\section{Scope of the arithmetic modular-form conclusion}\label{scope}
The modular-form assertion just proved uses the existing algebraic
metaplectic definition and an explicitly constructed finite level.
Its proof contains no appeal to complex realization of a cover,
upper-half-plane simple connectivity, analytic theta transformation,
or comparison of an analytic Fourier frame with a formal Tate frame.
Every required frame and every coefficient was instead constructed
inside one algebraic formal completion.

The conventional description of such objects by functions on the
upper half-plane belongs to the general comparison between arithmetic
and analytic definitions of modular forms. That comparison is not an
input to our coefficient theorem, our construction of the modular
forms, or the identification of their prescribed formal expansions.
In particular we do not present the comparison theorem itself as an
arithmetic theorem proved here. An analytic reformulation would use
that separate general dictionary; none is needed for the original
formal-series task formulated with arithmetic modular forms.

The finite frame level $Y$ may be noncongruence. Its finiteness and
etaleness do not imply domination by a principal congruence level,
and we have made no such assertion. Likewise, the specific Shimura
half-form lines in \cite[Definitions 3.4.1 and 3.4.4]{candelori}
include additional theta-multiplier data. We do not claim that the
common level has a prescribed Shimura multiplier. These distinctions
do not restrict the input family or the determination of all its
formal relations: the input question specifies neither a congruence
group nor a particular multiplier.

All computations accompanying this manuscript distinguish full-series
proofs from finite experiments. External PDFs, page locations, checksums,
and hypotheses are retained in \texttt{external\_sources/manifest.json}.
The prior reviewed manuscript and historical proof records are preserved;
no PASS label or finite test is used as a mathematical proof.
Lean formalization is deferred.

\begin{thebibliography}{9}
\bibitem{conrad} B. Conrad, \emph{Arithmetic moduli of generalized elliptic
curves}, Journal of the Institute of Mathematics of Jussieu
\textbf{6} (2007), 209--278. Sections 2.5 and 4.
\url{https://math.stanford.edu/~conrad/papers/genmoduli.pdf}.
\bibitem{kato} K. Kato, \emph{$p$-adic Hodge theory and values of zeta
functions of modular forms}, Ast\'erisque \textbf{295} (2004), 117--290.
Proposition 1.3 and Sections 1.4--1.10.
\url{https://numdam.org/item/AST_2004__295__117_0.pdf}.
\bibitem{mumford} D. Mumford, \emph{An analytic construction of degenerating
abelian varieties over complete rings}, Compositio Mathematica
\textbf{24} (1972), 239--272.
\url{https://numdam.org/item/CM_1972__24_3_239_0.pdf}.
\bibitem{klz} G. Kings, D. Loeffler and S. L. Zerbes,
\emph{Rankin--Eisenstein classes and explicit reciprocity laws},
arXiv:1503.02888, Theorem 4.1.1.
\url{https://arxiv.org/abs/1503.02888}.
\bibitem{candelori} L. Candelori, \emph{The transformation laws of algebraic
theta functions}, arXiv:1609.04486.
Remark 2.6.5 and Theorems 4.1.1, 4.1.3 supply the algebraic
full-level and projective comparison used in Section~\ref{refinement}.
The analytic comparison of Section 4.4 is not used in that section.
\url{https://arxiv.org/abs/1609.04486}.
\bibitem{candelori-key} L. Candelori, \emph{The algebraic functional
 equation of Riemann's theta function}, arXiv:1512.04415v3.
Section 2.2 fixes the genus-one even discriminant character.
\url{https://arxiv.org/abs/1512.04415}.
\end{thebibliography}
\end{document}
